\section{Open Directions}
\label{sec:open}

The present note provides a mathematical foundation for simplex memory networks. We outline two directions for future work: theoretical analysis of the architecture's fundamental properties, and design choices for practical deployment.

Theoretical questions concern the expressive power and continuum limits of $\SMN(n,m)$ architectures. An open problem is whether simplex memory networks satisfy a universal approximation theorem as $n,m \to \infty$, in analogy with classical results for feedforward networks \cite{Cybenko1989,HornikStinchcombeWhite1989}. The geometric constraints, namely fixed connectivity, potential-based orientation, and facet-attached interfaces, raise the question of whether these networks retain universal approximation capability. Another direction is to derive continuous-depth or mean-field limits, either via neural ODE formulations \cite{ChenRubanovaBettencourtDuvenaud2018} or transport/diffusion equations on simplicial domains \cite{EinizadeThanouMalliarosGiraldo2025}. Such limits would provide analytical tools for studying memory consolidation, decay, and long-time behavior.

Design choices concern architectural parameters and practical deployment. The scaling parameters $(n,m)$ control the interface-to-volume ratio $n/(n+m-1)$: large $m$ with small $n$ yields deep, narrow architectures, while small $m$ with large $n$ produces broad, shallow networks. Systematic exploration of this design space, guided by task complexity and memory requirements, would yield practical selection guidelines.
